\section{Cotangent extensions with zero ruling component}
\label{sec:elliptic-zero-ruling-extension}

Let $E=\mathbb C/\Lambda$, with $\Lambda=\mathbb Z+\mathbb Z\tau$ and $\operatorname{Im}\tau>0$. Let $A$ have degree $d>0$, and put $M=\mathbb P_{\mathrm{lines}}(\mathcal O_E\oplus A)$, with projection $f:M\to E$. Let $\eta,\theta$ be the class and form of \cref{lem:elliptic-additive-class}. Use the normalized factors of \cref{lem:elliptic-degree-automorphy-normalization}. Let $W$ be an extension of $\mathcal O_M$ by $\Omega_M^1$ with class $Kf^*\eta$, and put $V=W^*$. The homogeneous charts of \cref{thm:split-elliptic-cotangent-torsor}, now with zero ruling component, have transitions
\[
 v=1/t,\qquad u_\infty=u,\qquad
 y_\infty=-t^2y,\qquad w_\infty=w.
\]
For $\lambda=m\tau+n$, the lattice action is
\[
 (z,t,u,y,w)\longmapsto
 (z+\lambda,b_\lambda t,u+m(q_0ty+Kw),y/b_\lambda,w),
 \qquad q_0=2\pi\mathrm{i}d.
\]
For the local lift differences $z_i-z_j=m_{ij}\tau+n_{ij}$, the zero lifts in these charts differ by $Km_{ij}f^*\theta$ on base overlaps and agree on the two ruling charts, so they give the class $Kf^*\eta$. The linear parts are the cotangent transitions, and the inverse ruling chart gives the holomorphic lattice formula $u_\infty\mapsto u_\infty+m(Kw-q_0v y_\infty)$. Thus the same cocycle computation constructs $W$ also when its ruling component is zero. The fiberwise splitting used for a nonzero ruling component is unavailable, but ordinary cotangent polynomials give the homogeneous algebra directly.

\begin{lemma}\label{lem:elliptic-zero-ruling-quadric}
The graded algebra of fiber-homogeneous polynomial functions on $W$ over $\mathbb C\times\mathbb P^1$ is
\[
 \mathcal O(\mathbb C)[c_-,c_+,h,u,w]/(c_-c_+-h^2),
 \qquad c_-=y,\quad c_+=t^2y,\quad h=ty,
\]
where all five variables have degree one. The lattice fixes $h,w$, sends $c_-$ to $b_\lambda^{-1}c_-$ and $c_+$ to $b_\lambda c_+$, and translates $u$ by $m(q_0h+Kw)$.
\end{lemma}
\begin{proof}
Holomorphy in the infinite ruling chart bounds the degree in $t$ of the coefficient of $y^j$ by $2j$, as in \cref{lem:split-elliptic-tensor-weights}; the extra variable $w$ is unchanged in that chart. Every monomial $t^ay^j$ with $0\leq a\leq2j$ is a product of $y,ty,t^2y$, by distributing the exponent $a$ into $j$ integers in $\{0,1,2\}$. These generators satisfy $y(t^2y)=(ty)^2$.

In degree $j$ their function space has dimension $2j+1$. The polynomial algebra in three variables modulo the displayed nonzero quadratic relation has the same dimension
\[
 \binom{j+2}{2}-\binom{j}{2}=2j+1,
\]
where the second term is zero for $j<2$. The surjective graded map is consequently an isomorphism. Adjoining the two trivial fiber coordinates $u,w$ gives the stated algebra, with entire coefficients in $z$. In the other chart,
\[
 c_-=-v^2y_\infty,\qquad c_+=-y_\infty,
 \qquad h=-v y_\infty,
\]
so all generators are holomorphic. Substitution gives their lattice transformations.
\end{proof}

\begin{proposition}\label{prop:elliptic-zero-ruling-section-ring}
There is a degree-preserving algebra isomorphism
\[
 \bigoplus_{r\geq0}H^0(M,\operatorname{Sym}^rV)
 \simeq
 \left(\bigoplus_{r\geq0}H^0(M,\operatorname{Sym}^rT_M)\right)[w],
 \qquad \deg w=1.
\]
Consequently the dimension, volume and nonfinite-generation conclusions of \cref{cor:elliptic-extension-symmetric-dimensions} hold also for this extension.
\end{proposition}
\begin{proof}
Give $c_-,c_+$ weights $-1,1$, and $h,u,w$ weight zero. Reducing $c_-c_+$ by $h^2$ yields the negative blocks $c_-^kF(z,h,u,w)$, the positive blocks $c_+^kF(z,h,u,w)$ and the zero block $F$. They are unique: substituting $c_+=h^2/c_-$ embeds the relation algebra into $\mathcal O(\mathbb C)[h,u,w,c_-,c_-^{-1}]$, and distinct reduced weights have distinct Laurent powers with nonzero polynomial coefficients. In tensor degree $r$, the remaining polynomial degree is $r-k$.

The coefficient bundle $U_3$ has inverse translation $u\mapsto u-m(q_0h+Kw)$. Thus the weight $-k$ block is $H^0(E,A^k\otimes\operatorname{Sym}^{r-k}U_3)$, the weight $k$ block has $A^{-k}$ instead, and the zero block is $H^0(E,\operatorname{Sym}^rU_3)$. Put $H=h+(K/q_0)w$. Since $q_0\ne0$, this identifies $U_3$ with the previously defined $U$, using $H$ in place of its Euler variable, plus the trivial line generated by $w$.

By \cref{lem:elliptic-unipotent-symmetric-cohomology}, all positive weights vanish and the zero-weight algebra is $\mathbb C[H,w]$. Expanding the negative blocks in powers of $w$ identifies them with the tangent blocks of \cref{lem:split-elliptic-tensor-weights}. The substitutions $h\mapsto H$ and $y^k\mapsto c_-^k$ give bijections in every degree and surviving weight. Their products multiply the coefficient polynomials and add the powers of $c_-$, so these bijections preserve multiplication. This proves the algebra isomorphism. Summing the tangent dimensions and taking the quotient by $w$ gives the remaining assertions exactly as in \cref{cor:elliptic-extension-symmetric-dimensions}.
\end{proof}

\begin{proposition}\label{prop:elliptic-zero-ruling-torsor-cover}
The cotangent torsor $T_W$ is not Stein. Its universal cover is
\[
 \mathbb C^2\times\operatorname{Tot}
                      (\mathcal O_{\mathbb P^1}(-2)),
\]
which is also not Stein.
\end{proposition}
\begin{proof}
Set $w=1$. The subset $y=0$ in the finite ruling chart agrees with $y_\infty=0$ in the other chart. It descends to a submanifold fibered over
\[
 B_K=\mathbb C^2_{z,u}/\Lambda,\qquad
 (m\tau+n)\cdot(z,u)=(z+m\tau+n,u+Km),
\]
with fibers $\mathbb P^1$. The quotient is a complex manifold because translation in $z$ gives a free properly discontinuous action. Every ruling fiber is a compact connected curve in $T_W$, and holomorphic functions are constant on it. Thus $T_W$ is not holomorphically separable and is not Stein.

Before the lattice quotient, the affine charts have no affine ruling term and glue to $\mathbb C^2_{z,u}\times\operatorname{Tot}(\mathcal O_{\mathbb P^1}(-2))$. This space is simply connected: contraction of each vector-bundle fiber retracts its last factor onto $\mathbb P^1$, while $\mathbb C^2$ is contractible. The lattice action is free and properly discontinuous, so this is the universal cover. The zero section contains a compact projective line, again obstructing holomorphic separation and hence Steinness.
\end{proof}

Thus the zero and nonzero ruling classes have isomorphic graded section algebras, although \cref{thm:split-elliptic-cotangent-torsor} gives a Stein universal cover only for the nonzero ruling classes.
